\section{Connect validation, replay and publication}
Build the complete program in layers whose intermediate state is inspectable.
First admit the file header, then scan complete record envelopes and CRCs. Validate
all payload command streams and construct the directory without changing the display.
Only an accepted file should publish its frame count, period, duration and directory
count. A short directory destination must fail before any directory publication.

Next implement a record-application routine. For a full frame, reconstruct 4096
native display words from triples. For a delta, validate the command sequence, then
apply its already bounded changes to the previous working frame. Advance sequence
state only after successful application. Test a bad late command against a previously
successful image, retaining its actual pixels and sequence rather than resetting them
to a sentinel before the failure.

Add seek by loading the chosen full frame and replaying through the target. Compare
its exact words with sequential playback for every target in every public fixture.
The index construction cost belongs to opening the file; replay cost belongs to the
particular seek. Keep the file immutable across those operations, or rebuild the
validated directory when the source changes.

Finally connect the encoder and the presentation loop. Decode the encoder's result
and compare every RGB byte with its input frames. A valid round trip alone does not
prove the decoder rejects malformed files, so retain independent negative tests.
Then compare the scalar and MSA delta stages on identical inputs, followed by
whole-program size and work measurements. Construction through step 122 connects complete-file encoding to validation,
seeking, bounded storage, reload validation and timed sequential playback. Final
feedback, student/instructor distributions and manuscript integration remain.
